\scanpage{52}
% archive foliation: 51

À l'infini, la situation n'est pas très belle encore. Il vaut mieux dire:

Pour $\forall\,\xi_p\in\mathcal T_p$,
\deleted{la valeur en \unsure{$\infty$} de $\unsure{\mathfrak g}(\xi_p)$ est}
\inserted{Considérons la gerbe $h(\xi_p)$}
de groupe $\mu_2$ sur $\mathbb Z[1/p]$ définie par $(C(\xi_p),T(\xi_p))$,
\deleted{\unsure{Et soit}}\inserted{Considérons}
$h(\xi_p)_{\mathbb R}$, qui est une gerbe \uline{canonique} sur $\mathbb R$ de groupe $\mu_2$.
\deleted{J'ignore si}
\inserted{pour $\xi_p=$ élément de Deuring,}
elle a une description indépendante de la théorie des catégories de Deuring (et de celle des courbes elliptiques).
\unsure{Ceci nous dit} en tous cas qu'elle est non triviale; elle est, pour $\xi_p=$ élément de Deuring, munie d'une \unsure{plexion} $\xi_\infty$ (ce qui revient à donnée d'une classe d'isomorphisme \inserted{d'équivalences} avec la gerbe type $g_\infty$, \uline{mais il reste l'ambiguïté des signes}.

Catégorie
\[
C(\xi_p)_\infty\simeq
\text{catégorie des vectoriels \inserted{à droite $V$ de dim $1$} \deleted{\unsure{sur}} sur}
\]
$H$ ($=$ quaternions réels), avec \deleted{\unsure{hom}}\ $\nu$ foncteur norme réduite vers cat. de Picard de $\mathbb R$.

Ainsi il y a un foncteur ``canonique au signe près''
\diagram{\begin{tikzcd}[column sep=large]
C(\xi_p) \arrow[r,"\xi_\infty^C"] &
\deleted{\unsure{C(\mathbb R)}\ \text{catégorie}}\ \unsure{\mathcal P(\mathbb R)|\mathcal Q(\mathbb R)}\\
\deleted{T(\xi_p)} \arrow[r,mapsto] &
\deleted{\unsure{\mathbb R}}
\end{tikzcd}}
\margin{\unsure{n'est pas spécial} à $C(\xi_p)$, \deleted{\unsure{vaut}} pour tt $C$.}
\margin{la factorisation \unsure{passe} par les vect. quaternioniens de dim $1$ \uuline{unitaires}}
\deleted{\unsure{qui est équivalente à celle des} \unsure{$\mathbb R$} \unsure{définit un autre} \unsure{et isom}}
\unsure{tout autre choix} \unsure{à ces}, l'isom \uline{défini au signe près} (chic!).

Pour $\xi_p=$ élément de \inserted{Deuring}\deleted{\unsure{Frobenius}}, on obtient un foncteur \inserted{ess. canonique} des courbes elliptiques de Hasse nul vers les vectoriels \inserted{\unsure{quatern. unitaires}} de dim $1$.
La \deleted{forme} Norme correspond à la polarisation, la transposition, la transposition
